A inflação do balão de raio $1$ para raio $R$ é uma transformação de escala uniforme (homotetia), em que cada ponto da superfície é afastado do centro por um fator multiplicativo $R$.

\textbf{Prova de invariância do ângulo:}
Sejam $\mathbf{t}_1$ e $\mathbf{t}_2$ os vetores tangentes a duas curvas que se cruzam em um ponto da superfície antes da inflação. O ângulo $\theta$ entre eles é:
\[
\cos\theta = \frac{\mathbf{t}_1\cdot\mathbf{t}_2}{|\mathbf{t}_1|\,|\mathbf{t}_2|}.
\]
Após escalar de raio $1$ para $R$, todos os vetores tangentes são multiplicados pelo mesmo fator $R$:
\[
\mathbf{t}_i \;\to\; R\,\mathbf{t}_i.
\]
O novo ângulo $\theta'$ satisfaz:
\[
\cos\theta' = \frac{(R\mathbf{t}_1)\cdot(R\mathbf{t}_2)}{|R\mathbf{t}_1|\,|R\mathbf{t}_2|} = \frac{R^2\,(\mathbf{t}_1\cdot\mathbf{t}_2)}{R^2\,|\mathbf{t}_1|\,|\mathbf{t}_2|} = \cos\theta.
\]
Portanto, $\theta' = \theta$: o \textbf{ângulo inscrito permanece invarídvel}.

As grandezas que mudam com a inflação são:
\begin{itemize}
\item Comprimentos: escalam por $R$;
\item Áreas: escalam por $R^2$.
\end{itemize}
Ângulos são razões adimensionais (arco/raio) e são preservados por qualquer escala uniforme.

